@ID: ON16D1P1U
@BIDANG: Aljabar Linear

Misalkan $A\in\mathbb R^{2\times2}$ dan $b\in\mathbb R^2$, $b\ne0$. Pemetaan $T:\mathbb R^2\to\mathbb R^2$ didefinisikan sebagai

$$
T(x)=Ax+b.
$$

Oleh $T$, setiap titik dipetakan ke titik lain dan setiap garis dipetakan ke garis lain. Buktikan bahwa

$$
Ab\ne b
$$

dan

$$
(A-I)^2=0,
$$

di mana $I$ adalah matriks identitas $2\times2$.

@ID: ON16D1P2U
@BIDANG: Struktur Aljabar

Misalkan $R$ suatu ring. Unsur $x\ne0$ di $R$ disebut sebagai pembagi nol kiri jika terdapat $a\ne0$ di $R$ sedemikian sehingga $xa=0$. Secara serupa, unsur $y\ne0$ di $R$ disebut sebagai pembagi nol kanan jika terdapat $b\ne0$ di $R$ sedemikian sehingga $by=0$.

Misalkan $x$ pembagi nol kiri dan $y$ pembagi nol kanan di $R$. Jika $R$ gelanggang dengan berhingga banyak unsur dan $xy\ne0$, tunjukkan bahwa $xy$ sekaligus merupakan pembagi nol kiri dan pembagi nol kanan.

@ID: ON16D1P3U
@BIDANG: Analisis Real

Misalkan fungsi kontinu $f:[0,1]\to\mathbb R^+$. Definisikan

$$
u_n=\left(\int_0^1(f(x))^n\,dx\right)^{1/n}
$$

dan

$$
M=\sup_{0\le x\le1}f(x).
$$

Buktikan bahwa:

(a) Untuk setiap $\varepsilon>0$, terdapat bilangan $\alpha$ dan $\beta$ dengan $0\le\alpha<\beta\le1$ sedemikian sehingga

$$
M(1-\varepsilon)\le f(x)\le M,\qquad \forall x\in(\alpha,\beta).
$$

(b)

$$
\lim_{n\to\infty}u_n=M.
$$

@ID: ON16D1P4U
@BIDANG: Analisis Real

Untuk setiap pernyataan berikut, buktikan atau berikan sebuah contoh penyangkal.

(a) Jika fungsi kompleks $f(z)$ terdefinisi pada cakram satuan terbuka

$$
D=\{z\in\mathbb C:|z|<1\}
$$

dan jika $f^2(z)$ analitik pada $D$, maka $f$ analitik pada $D$.

(b) Jika fungsi kompleks $f$ terdefinisi dan terdiferensial kontinu pada cakram satuan terbuka

$$
D=\{z\in\mathbb C:|z|<1\}
$$

dan jika $f^2(z)$ analitik pada $D$, maka $f$ analitik pada $D$.

@ID: ON16D1P5U
@BIDANG: Kombinatorika

Setiap bilangan $1,2,\ldots,n$ diwarnai dengan tiga warna. Tentukan $n$ terkecil sedemikian sehingga untuk sebarang pewarnaan dari bilangan-bilangan tersebut, selalu terdapat tiga bilangan $x_1,x_2,x_3\in\{1,2,\ldots,n\}$ yang berwarna sama dan memenuhi

$$
x_1+x_2=x_3.
$$

@ID: ON16D2P1U
@BIDANG: Aljabar Linear

Misalkan $A=[a_{ij}]\in\mathbb R^{n\times n}$ memiliki sifat terdapat bilangan-bilangan real konstan $c_1,c_2,\ldots,c_n$ sehingga

$$
\sum_{i=1}^{k}a_{ij}=c_k,
\qquad
\text{untuk semua }1\le j\le k\le n.
$$

Tentukan semua nilai eigen dari $A$.

@ID: ON16D2P2U
@BIDANG: Struktur Aljabar

Misalkan $G$ suatu grup dengan $n$ unsur ($n\ge2$) dan $p$ adalah faktor prima terkecil dari $n$. Misalkan $G$ memiliki tepat satu subgrup $H$ yang berorde $p$. Tunjukkan bahwa untuk setiap $h\in H$ dan $g\in G$ berlaku

$$
hg=gh.
$$

@ID: ON16D2P3U
@BIDANG: Analisis Real

Diberikan suatu konstanta positif $K$. Misalkan barisan bilangan real nonnegatif $\{a_n\}$ memenuhi

$$
a_{m+n}\le a_m+a_n+K,
$$

untuk setiap bilangan positif $m,n$. Perlihatkan bahwa barisan $\left\{\frac{a_n}{n}\right\}$ konvergen.

@ID: ON16D2P4U
@BIDANG: Analisis Kompleks

Diberikan bilangan-bilangan kompleks $z_1,z_2,z_3$ dengan $z_1\ne z_2$ dan $|z_1|=|z_2|=|z_3|=r>0$. Buktikan bahwa

$$
\min\{|kz_1+(1-k)z_2-z_3|:k\in\mathbb R\}
=
\frac{|z_1-z_3||z_2-z_3|}{2r}.
$$

@ID: ON16D2P5U
@BIDANG: Kombinatorika

Misalkan $G=(V(G),E(G))$ adalah suatu graf terhubung tak trivial. Didefinisikan $c:E(G)\to\{1,2,\ldots,k\}$, $k\in\mathbb N$, sebagai suatu pewarnaan sisi di mana dua sisi yang bertetangga boleh berwarna sama.

Suatu lintasan $P$ dari titik $u$ ke titik $v$ ditulis sebagai lintasan $u-v$ di $G$. Lintasan tersebut dinamakan *rainbow path* jika tidak terdapat dua sisi di $P$ yang berwarna sama.

Graf $G$ disebut *rainbow connected* jika setiap dua titik yang berbeda di $G$ dihubungkan oleh *rainbow path*. Pewarnaan sisi yang mengakibatkan $G$ bersifat *rainbow connected* dikatakan *rainbow coloring*. Bilangan *rainbow connection* dari graf terhubung $G$, ditulis $rc(G)$, yaitu banyaknya warna minimal yang diperlukan agar graf $G$ menjadi *rainbow connected*.

Misalkan $c$ adalah suatu *rainbow coloring* dari graf terhubung $G$. Misalkan $C_n$ adalah graf lingkaran (*cycle*) dengan $n\ge3$. Tentukan $rc(C_n)$, kemudian buktikan.